\chapter{BV coarea and level-set calculus}\label{sec:bv-coarea}
Throughout this section $n\geq2$. A BV function packages a family of finite-perimeter superlevel sets. The coarea identity makes this precise and turns approximation of functions into approximation of interfaces. We follow \cite{Giu84}, Theorems 1.23--1.24, pp.~20--23, with the superlevel convention $E_t=\{u>t\}$ instead of the source's sublevel convention. Perimeter does not require the set to have finite volume.

\begin{lemma}[Layer cake and the distance between level sets]\label{lem:level-set-layer-cake}
For real numbers $a,b$,
\[
 a=\int_{\mathbb R}(\chi_{\{a>t\}}-\chi_{\{0>t\}})\,dt,
 \qquad |a-b|=\int_{\mathbb R}|\chi_{\{a>t\}}-\chi_{\{b>t\}}|\,dt.
\]
Consequently, for measurable finite-valued functions $u,v$ on an open set $W$,
\[
 \|u-v\|_{L^1(W)}
 =\int_{\mathbb R}|(\{u>t\}\mathbin{\triangle}\{v>t\})\cap W|\,dt,
\]
including infinite values.
\end{lemma}
\begin{proof}
The indicators in the first formula differ only on the interval between $0$ and $a$, with the appropriate sign. The absolute difference in the second is the indicator of the interval between $a$ and $b$, up to its endpoints. Integrate the pointwise identity using Tonelli. This is the level-set identity used in \cite{Giu84}, the proof of Theorem 1.23, pp.~20--21.
\end{proof}

\begin{proposition}[Coarea for smooth superlevel sets]\label{prop:smooth-superlevel-coarea}
If $u\in C^\infty(U)$ and $W\subset U$ is open, then
\[
 \int_W|\nabla u|\,dx=\int_{\mathbb R}P(\{u>t\};W)\,dt.
\]
For every regular value $t$, the superlevel boundary inside $U$ is the smooth level hypersurface, and its outward normal is $-\nabla u/|\nabla u|$.
\end{proposition}
\begin{proof}
At a regular value, the implicit function theorem and the local divergence theorem give
$P(\{u>t\};W)=\mathcal H^{n-1}(W\cap\{u=t\})$. The outward normal points toward decreasing $u$. Sard's theorem excludes the nonregular values up to a null set. Apply Theorem \ref{thm:coarea-formula} with $f=u$, whose coarea Jacobian is $|\nabla u|$. Exhaustion handles noncompact $W$ and possibly infinite integrals. Compare the smooth step in \cite{Giu84}, Theorem 1.23.
\end{proof}

\begin{lemma}[Almost-everywhere convergence of superlevel sets]\label{lem:level-set-convergence}
If $u_j\to u$ in $L^1(U)$, there is a subsequence such that
\[
 |(\{u_j>t\}\mathbin{\triangle}\{u>t\})\cap U|\longrightarrow0
 \quad\text{for almost every }t\in\mathbb R.
\]
The same conclusion holds locally when the convergence is in $L^1_{\mathrm{loc}}(U)$, with a single exceptional null set for a countable exhaustion.
\end{lemma}
\begin{proof}
Choose a subsequence with $\sum_j\|u_j-u\|_{L^1(U)}<\infty$. The layer-cake identity and Tonelli show that for almost every $t$ the sum of the displayed symmetric-difference measures is finite; hence its terms tend to zero. For local convergence choose a diagonal subsequence whose errors are summable on each set of a countable relatively compact exhaustion. The union of its exceptional sets is still null. See \cite{Giu84}, p.~21.
\end{proof}

\begin{theorem}[Whole-space BV coarea]\label{thm:bv-coarea-global}
For $u\in BV(\mathbb R^n)$,
\[
 |Du|(\mathbb R^n)=\int_{\mathbb R}P(E_t;\mathbb R^n)\,dt,
 \qquad E_t=\{u>t\}.
\]
In particular almost every $E_t$ has finite perimeter. Its volume may be infinite, especially when $t<0$.
\end{theorem}
\begin{proof}
For a compactly supported smooth vector field $X$, the signed layer-cake representation gives
\[
 \int u\operatorname{div}X
 =\int_{\mathbb R}\int_{E_t}\operatorname{div}X\,dx\,dt,
\]
because the baseline term integrates to zero. Fubini is justified before this cancellation by the bound $|u||\operatorname{div}X|$ on the compact support. If $|X|\leq1$, the right side is at most $\int P(E_t)dt$. Taking the supremum proves one inequality. Measurability of $t\mapsto P(E_t;W)$ for an open $W$ follows by taking the supremum over a countable dense family of compactly supported $C^1$ test fields in the variational definition.

For the reverse inequality, let $u_j$ be smooth mollifications converging strictly as in Corollary \ref{cor:bv-strict-mollification}. By Lemma \ref{lem:level-set-convergence}, pass to a subsequence whose superlevel sets converge in symmetric difference for almost every $t$. Perimeter lower semicontinuity, Fatou, and smooth coarea give
\[
 \int P(E_t)\,dt\leq\liminf_j\int P(\{u_j>t\})\,dt
 =\lim_j\int|\nabla u_j|=|Du|(\mathbb R^n).
\]
This proves the identity. The argument is the approximation proof of \cite{Giu84}, Theorem 1.23, with strict approximation supplied here.
\end{proof}

\begin{theorem}[Local and weighted BV coarea]\label{thm:bv-coarea-local}
If $u\in BV_{\mathrm{loc}}(U)$, then almost every $E_t=\{u>t\}$ has locally finite perimeter in $U$, and for every Borel $A\subset U$,
\[
 |Du|(A)=\int_{\mathbb R}|D\chi_{E_t}|(A)\,dt.
\]
Equivalently, for every nonnegative Borel function $g$,
\[
 \int_Ug\,d|Du|
 =\int_{\mathbb R}\int_Ug\,d|D\chi_{E_t}|\,dt.
\]
The exceptional levels can be fixed once for a countable exhaustion of $U$.
\end{theorem}
\begin{proof}
First take $u\in BV(\mathbb R^n)$. The previous theorem makes
$\nu(A)=\int|D\chi_{E_t}|(A)dt$ a finite Borel measure. To justify this measure construction, measurability is known for open $A$ by the countable test-field description; regular approximation and the monotone-class theorem extend it to Borel sets, and Tonelli gives countable additivity. Values at the exceptional levels are set to zero.

For an open ball $B$ with $|Du|(\partial B)=0$, apply the layer-cake test argument with fields supported in $B$ to obtain $|Du|(B)\leq\nu(B)$. The mollification and Fatou argument localized to $B$ gives
\[
 \nu(B)\leq\liminf_j\int_B|\nabla u_j|
 \leq\lim_{\varepsilon\downarrow0}|Du|(B_\varepsilon)
 =|Du|(\overline B)=|Du|(B).
\]
Here $B_\varepsilon$ is the $\varepsilon$ neighborhood and the middle estimate is Lemma \ref{lem:bv-mollification}. At each center only countably many radii have positive boundary mass. Approximation from below by the other radii gives equality on every ball; finite Radon measures agreeing on all balls agree on Borel sets. Thus $\nu=|Du|$.

For local $u$, choose a compactly supported smooth cutoff equal to one near any prescribed relatively compact open set. Proposition \ref{prop:bv-cutoff-product} gives a whole-space BV function agreeing with $u$ there. Its level sets and derivative measures therefore agree locally with those of $u$. A countable exhaustion proves the stated identity and the common exceptional-set assertion throughout $U$. Increasing simple functions and monotone convergence give the weighted version. This local form refines \cite{Giu84}, Theorem 1.23.
\end{proof}

\begin{proposition}[Lipschitz composition preserves local BV]\label{prop:bv-lipschitz-composition}
Let $u\in BV_{\mathrm{loc}}(U)$ and let $F:\mathbb R\to\mathbb R$ have Lipschitz constant $L$. Then $F(u)\in BV_{\mathrm{loc}}(U)$ and
\[
 |D(F(u))|\leq L|Du|
\]
as measures. If $u\in BV(U)$ and $F(0)=0$, then $F(u)\in BV(U)$ as well.
\end{proposition}
\begin{proof}
First let $F$ be smooth with $|F'|\leq L$. Mollify a cutoff of $u$ that agrees with it near a ball $B\Subset U$ with $|Du|(\partial B)=0$. The smooth chain rule gives $|\nabla F(u_j)|\leq L|\nabla u_j|$, and Lipschitz continuity gives $F(u_j)\to F(u)$ in $L^1(B)$. Lower semicontinuity and the localized mollification bound yield $|D(F(u))|(B)\leq L|Du|(B)$. Approximate radii to obtain this bound on every relatively compact ball. Integrating these ball inequalities against a nonnegative smooth function of the center, dividing by the ball volume, and shrinking the radius yields the inequality on smooth test functions. Radon measure regularity gives the measure inequality.

For general $F$, convolve it on the real line. The smooth functions retain Lipschitz constant $L$ and converge uniformly to $F$; another lower-semicontinuity argument proves the result locally. Finally $|F(u)|\leq L|u|$ when $F(0)=0$, giving global integrability and the global variation bound. This supplies the composition step used with the coarea calculus of \cite{Giu84}, Chapter 1; it does not claim the full BV chain rule across jumps.
\end{proof}

\begin{corollary}[Exact variation of a truncation]\label{cor:bv-truncation-coarea}
For finite $a<b$, put $T_{a,b}(s)=\min\{\max\{s-a,0\},b-a\}$. If $u\in BV_{\mathrm{loc}}(U)$, then for Borel $A\subset U$,
\[
 |D(T_{a,b}(u))|(A)=\int_a^bP(\{u>t\};A)\,dt\leq|Du|(A).
\]
Here $P(E;A)$ denotes the perimeter measure $|D\chi_E|(A)$ for almost every admissible level.
\end{corollary}
\begin{proof}
Composition is locally BV by Proposition \ref{prop:bv-lipschitz-composition}. For $0<s<b-a$, its superlevel set is $\{u>a+s\}$; outside this interval the superlevel set is either all of $U$ or empty, up to the endpoint levels. These two sets have zero relative perimeter. Apply local coarea and change variables $t=a+s$.
\end{proof}

\begin{corollary}[Positive and negative variations split exactly]\label{cor:bv-positive-negative}
For $u\in BV_{\mathrm{loc}}(U)$, write $u^+=\max\{u,0\}$ and $u^-=\max\{-u,0\}$. Then
\[
 |Du|=|Du^+|+|Du^-|
\]
as measures; in particular $|D|u||\leq|Du|$.
\end{corollary}
\begin{proof}
The nonnegative part has the same positive superlevel sets as $u$, so coarea gives $|Du^+|(A)=\int_0^\infty P(\{u>t\};A)dt$. For the negative part use $\{u<-s\}$ and complementation. Except for a countable set of levels, $\{u=t\}$ is Lebesgue-null locally: on each finite-volume exhaustion set there can be only countably many disjoint level sets of positive measure. Thus $\{u<t\}$ and the complement of $\{u>t\}$ agree almost everywhere for almost every $t$ and have the same perimeter. This gives the negative half of the coarea integral. Summing proves the identity. Since $|u|=u^++u^-$, the triangle inequality for vector measures gives the last conclusion.
\end{proof}

\begin{proposition}[The BV lattice inequality]\label{prop:bv-lattice}
If $u,v\in BV_{\mathrm{loc}}(U)$, then their maximum and minimum are locally BV and
\[
 |D\max\{u,v\}|+|D\min\{u,v\}|\leq|Du|+|Dv|
\]
as measures.
\end{proposition}
\begin{proof}
For smooth $u,v$, approximate the maximum and minimum by
\[
 M_\delta=\frac{u+v+\sqrt{(u-v)^2+\delta^2}}2,
 \qquad m_\delta=\frac{u+v-\sqrt{(u-v)^2+\delta^2}}2.
\]
Their gradients are complementary convex combinations of $\nabla u,\nabla v$, so
$|\nabla M_\delta|+|\nabla m_\delta|\leq|\nabla u|+|\nabla v|$. Let $\delta\downarrow0$ using local $L^1$ convergence and lower semicontinuity. For BV inputs, mollify cutoffs agreeing with $u,v$ near a ball whose boundary has zero $(|Du|+|Dv|)$ measure. The maxima and minima converge in $L^1$, because these operations are Lipschitz in their two scalar arguments. The same localized mollification estimate as above proves the bound on that ball. Approximate radii, integrate the ball inequalities against nonnegative smooth functions of their centers, and shrink the radius. This gives the measure inequality on smooth test functions and hence on Borel sets by regularity. This is an approximation consequence of the BV framework in \cite{Giu84}, Chapter 1.
\end{proof}

\begin{corollary}[Union and intersection of finite-perimeter sets]\label{cor:perimeter-submodularity}
If $E,F$ have locally finite perimeter in $U$, then so do $E\cup F$ and $E\cap F$, and
\[
 P(E\cup F;A)+P(E\cap F;A)\leq P(E;A)+P(F;A)
\]
for every Borel $A\subset U$.
\end{corollary}
\begin{proof}
The characteristic functions of the union and intersection are respectively the maximum and minimum of $\chi_E,\chi_F$. Apply Proposition \ref{prop:bv-lattice}. This argument is insensitive to changes on Lebesgue-null sets.
\end{proof}

\begin{theorem}[Smooth approximation of a bounded finite-perimeter set]\label{thm:perimeter-smooth-approximation}
Let $E\subset\mathbb R^n$ be bounded up to a null set and have finite total perimeter. There are bounded open sets $E_j$ with smooth compact boundary such that
\[
 |E_j\mathbin{\triangle}E|\to0,\qquad P(E_j;\mathbb R^n)\to P(E;\mathbb R^n).
\]
This is \cite{Giu84}, Theorem 1.24, pp.~22--23. The conclusion does not prescribe the topology of the approximating boundaries.
\end{theorem}
\begin{proof}
Let $v_\varepsilon=\rho_\varepsilon*\chi_E$. It is smooth, compactly supported, and takes values in $[0,1]$. Choose $\delta_j\downarrow0$ with $\delta_j<1/2$ and then $\varepsilon_j\downarrow0$ so that $\|v_{\varepsilon_j}-\chi_E\|_1/\delta_j\to0$. Smooth coarea and the mollification bound give
\[
 \int_{\delta_j}^{1-\delta_j}P(\{v_{\varepsilon_j}>t\})dt\leq P(E).
\]
There is a regular value $t_j$ in this interval with perimeter at most $P(E)/(1-2\delta_j)+1/j$. Set $E_j=\{v_{\varepsilon_j}>t_j\}$. Sard and the implicit function theorem give its smooth compact boundary. Pointwise disagreement of its characteristic function with $\chi_E$ forces $|v_{\varepsilon_j}-\chi_E|\geq\delta_j$. Hence
$|E_j\triangle E|\leq\delta_j^{-1}\|v_{\varepsilon_j}-\chi_E\|_1\to0$.
The selected upper bound and perimeter lower semicontinuity now give convergence of perimeters. The chosen levels may vary with $j$; no convergence claim at a preassigned threshold is needed.
\end{proof}

\begin{example}[A jump contributes its height times perimeter]\label{ex:bv-jump-coarea}
Let $E$ have locally finite perimeter in $U$ and let $a>b$. For $u=b+(a-b)\chi_E$, the superlevel set is $U$, $E\cap U$, or empty according as $t<b$, $b<t<a$, or $t>a$. Coarea gives $|Du|=(a-b)|D\chi_E|$, in agreement with distributional differentiation. A jump is recorded by the perimeter measures of an interval of levels, not by a classical gradient almost everywhere in the ambient domain.
\end{example}

\begin{example}[A fixed level can fail under strong convergence]\label{ex:exceptional-level}
On a bounded nonempty open $U$, the constants $u_j=1/j$ converge to $u=0$ in $L^1(U)$ and strictly in BV. At level $t=0$, however, $\{u_j>0\}=U$ and $\{u>0\}$ is empty. Thus even strict convergence does not force convergence of every superlevel set. The almost-everywhere qualification in Lemma \ref{lem:level-set-convergence} is necessary.
\end{example}
